% page 487 of the facsimile (image p-486 of the scan)
% block 167.6 x 233.3 mm ; left margin 34.4 mm
\pgc{12.700mm}{0.000mm}{
\l{\marge[76.860mm]{36.492mm}{\ornamento{12.252mm}{ornaments/letters/letter-R.png}}}
\l{}
\l{}
\l{}
\l{}
\l{}
\l{}
\l{}
\l{}
\l{}
\l{}
\l{}
\l{\sou{rabata}r. (trans., ek, de).(\sou{komerco}) (\sou{aludante furnisisto}).Grantar}
\l{\cel{3}preco-diminuto a sua kompranto. Grantar konvencione (a kom-}
\l{\cel{3}pranto) diminuto di la pekunio-quanto debata, kondicione ke ica}
\l{\cel{3}pagesos ante la dato normala. - DEFIRS.}
\l{}
\l{\sou{rabio}. (\sou{patol.}) Morbo virulenta qua aparas spontane che la hundo}
\l{\cel{3}(plu rare che la kato, la volfo, e c.), komunikebla per mordo}
\l{\cel{3}da ta animali, e finas ye acesi hororinda di konvulsi, e ye}
\l{\cel{3}morto. (Morbo vinkita da Louis Pasteur).- EFIS.}
\l{}
\l{\sou{rabino.}\cel{2}I. Doktoro pri la lego religiala di la Judi. - II. Dok-}
\l{\cel{3}toro pri la kulto Judala, chefo di socio de Judi. - DEFIRS.}
\l{}
\l{\sou{raboto.}\cel{2}(\sou{tekn.)}\cel{2}Utensilo di menuzisto, cizelo ek stalo, muntita}
\l{\cel{3}oblique interne di ligno-peco rektangula kavigita, qua lasas}
\l{\cel{3}pasar la parto tranchiva ; uzata por planigar, diminutar la}
\l{\cel{3}areo di ligno-peco. - F.}
\l{}
\l{\sou{racemo}. I. (\sou{biol.}) Nomo ciencala di la grapo. - II. (\sou{kemio).}}
\l{\cel{3}Substanco qua ne efikas a la lumo polarigita, konsistanta ye}
\l{\cel{3}kombinuro (kun molekuli inter-egala) di du inversaji optikala.}
\l{\cel{3}DEFIS.}
\l{}
\l{\sou{raciono}. La povo dicernar la exaktaji : (a) la fakultato desko-}
\l{\cel{3}vrar lo exakta ; (b)(\sou{filoz.}) la fakultato konceptar la exaktaji}
\l{\cel{3}absoluta. - DEFIRS.}
\l{}
\l{\sou{rado}. Parto de maro,qua penetras aden la tero ferma, e formacas}
\l{\cel{3}quaza doko naturala ube la navi povas shirmar su. - DFIRS.}
\l{}
\l{\sou{radiar}. (\sou{netrans}.)I. (\sou{aludante lumo-sprici rekta}).Lansesar quale}
\l{\cel{3}dardo, de centro lumoza, diverge, omna-sinse. - II. (\sou{fiziko)}}
\l{\cel{3}Lansar kaloro-sprici. - II.(\sou{geom.})\cel{2}(\sou{pri rekti)}. Irar de la}
\l{\cel{3}centro di cirklo a la cirklo-kurvo, a la surfaco di la afero.}
\l{\cel{3}- DEFIRS.}
\l{}
\l{\sou{radiatoro}. (\sou{teknol}.) I. Organo di ensemblo mekanikala, destinita}
\l{\cel{3}ad eventigar la koldesko di liquido o di gaso uzata, en la en-}
\l{\cel{3}semblo, kom vehilo di kaloro. - II. Hem-aparato uzata por}
\l{\cel{3}varmigar la atmosfero di la chambri per fluido (aero, aquo,}
\l{\cel{3}vaporo) qua uzesas kom vehilo di kaloro por transportar olca}
\l{\cel{3}kelka-diste. -\cel{2}F.}
}
